\section{Liste Concatenate: Esercizi Avanzati (Lab 13)}

\subsection{Manipolazione di Liste}

\subsubsection{1. Elemento N-esimo (listNth)}
\begin{concept}{Accesso Sequenziale}
Trovare l'elemento all'indice `n` riducendo `n` ad ogni passo.
\end{concept}
\begin{jsbox}
function listNth(nodo, n) {
    if (!nodo) return null;
    if (n === 0) return nodo; // Arrivati!
    return listNth(nodo.next, n - 1);
}
\end{jsbox}

\subsubsection{2. Verifica Ordinamento (listIsSorted)}
\begin{concept}{Verifica Monotona}
Controllare se ogni nodo è minore del successivo.
\end{concept}
\begin{jsbox}
function listIsSorted(nodo) {
    if (!nodo || !nodo.next) return true; // 0 o 1 elemento: OK
    
    // Se il corrente è maggiore del prossimo -> Falso
    if (nodo.val > nodo.next.val) return false;
    
    // Controllo il resto
    return listIsSorted(nodo.next);
}
\end{jsbox}

\subsubsection{3. Mappare valori (listMap)}
\begin{concept}{Mappatura Ricorsiva}
Creare una NUOVA lista applicando $f(x)$ a ogni nodo.
\end{concept}
\begin{jsbox}
function listMap(nodo, f) {
    if (!nodo) return null;
    
    return {
        val: f(nodo.val),       // Trasformo il valore
        next: listMap(nodo.next, f) // Ricorro sul resto
    };
}
\end{jsbox}

\subsubsection{4. Rimuovere pezzi (listSplice)}
\begin{concept}{Taglio Chirurgico}
Rimuovere `count` nodi a partire dalla posizione `start`.
Ci sono 3 casi: Prima del taglio, Nel taglio, Fine.
\end{concept}
\begin{jsbox}
function listSplice(nodo, start, count) {
    if (!nodo) return null; // Fine lista
    
    // Caso 1: Devo ancora arrivare al punto di taglio
    if (start > 0) {
        nodo.next = listSplice(nodo.next, start - 1, count);
        return nodo;
    }
    
    // Caso 2: Sono nel punto di taglio, devo rimuovere
    if (count > 0) {
        // Salto 'nodo' e chiamo la funzione sul prossimo
        return listSplice(nodo.next, 0, count - 1);
    }
    
    // Caso 3: Finito di tagliare (count == 0)
    return nodo;
}
\end{jsbox}
